\section{Related work}
\label{sec:related}

\paragraph{Functional tensor decomposition.}
FunBaT~\citep{funbat2024} replaces the discrete factor tables of a Tucker model
with GP-distributed latent functions, so unobserved coordinates are interpolated
rather than looked up, and performs inference by message passing on an
equivalent state-space form.  RR-FBTC~\citep{rrfbtc2025} extends functional
Bayesian tensor completion with automatic rank determination.  Both use
\emph{generic} per-mode kernels; neither derives the kernels from the operator
that generated the field, which is the only thing this paper changes.  We hold
the host model, ranks, feature budget and optimiser fixed and vary only where
the spectra come from, so the comparison is against those methods' kernels
rather than against those methods.

\paragraph{GP priors that encode a PDE.}
EPGP~\citep{epgp2023} constructs a GP whose samples satisfy a constant-coefficient
linear PDE exactly, via the Ehrenpreis--Palamodov representation.  Tensor
GPs~\citep{tensorgp2025} solve nonlinear PDEs by collocation with
one-dimensional GP factors and a tensor decomposition.  Both require the equation
to be known completely, coefficients included, and both aim at solving rather
than at completing noisy sparse observations.  We ask for less and deliver less:
only the \emph{form} of the operator, and only an approximation to the
second-order spectrum of its solutions, which is what allows the coefficients to
be wrong by $50\%$ throughout our experiments.

\paragraph{Learning expressive stationary kernels.}
Spectral mixture kernels~\citep{spectralmixture2013} parameterise the spectral density
directly and are dense in the class of stationary kernels, so no fixed
construction can beat them on expressiveness -- only on sample efficiency, and
only where data is too scarce to identify the parameters.  We therefore include a
matched spectral-mixture dictionary as a baseline rather than claiming novelty
for using a spectral representation.  Our atoms are a nonnegative low-rank
projection of an operator family's joint spectrum, which is a physically
reachable set rather than a freely parameterised one.

\paragraph{Sensor placement and observability.}
Optimal sensor placement asks where to put sensors to make a field observable.
We ask the complementary question, which is what the applications force: the
placement is fixed and bad, and the question is what prior makes the resulting
extrapolation possible.  The closest point of contact is that our feasibility
screen is an observability test -- we reject fields whose boundary carries a
disproportionately small share of the energy, because no prior recovers what the
sensors never see.

\paragraph{What is not claimed.}
Neither ``GP plus tensor'' nor ``PDE-informed GP'' is new, and the title and
abstract do not claim them.  The narrow claim is the construction -- a
nonnegative, PSD-preserving low-rank projection of a non-separable joint
operator spectrum into per-mode kernels, usable when only the equation's form is
known -- together with the empirical finding about when that matters, which is
when the sensor geometry makes hyper-parameter selection impossible.
